\section{Planned Validation Study}
\label{sec:protocol}

The analyses above use public laboratory data and simulation. Whether the approach helps real people in daily life needs a prospective study. We describe the planned design here; before any data are collected it will be registered on OSF Registries with a full statistical analysis plan, under an institutional sponsor and with approval from a human research ethics committee, and only then will it be pre-specified.

\paragraph*{Intended use and population} BlinkTwin is intended for self-monitoring by adults during desk work at a laptop. The study recruits adults aged 18--60 who work at a laptop for at least 4~h a day. Exclusions: night-shift work or transmeridian travel within two weeks; diagnosed sleep disorders; a positive screen for obstructive sleep apnoea (STOP-Bang~\cite{chung2008stop}); excessive daytime sleepiness (Epworth Sleepiness Scale above 10~\cite{johns1991new}); clinically relevant insomnia (Insomnia Severity Index of 15 or more~\cite{bastien2001validation}); pregnancy; and medication that affects alertness. Chronotype is measured with the Morningness--Eveningness Questionnaire~\cite{horne1976self}. Claims will be limited to this population.

\paragraph*{Field phase} Each participant uses BlinkTwin on their own computer for 10 working days in \emph{study mode}, which records events but hides every estimate, forecast, check and suggestion, so that the display cannot anchor ratings or change behaviour. Wrist actigraphy and a sleep diary record sleep; caffeine, alcohol and naps are logged. On each day the participant runs the webcam during laptop work and completes four check-ins at fixed times (09:30, 12:00, 14:30 and 17:00) and two \emph{probe} ratings at random times between them. Probe ratings are withheld from the model and used only for scoring, so that predictions between check-ins, where passive signals can matter, are evaluated. At each check-in the participant rates the KSS first, then takes the 3-min PVT-B, then rates NASA-TLX for the work block just finished. Each device's display and input latency is measured once with a photodiode and a key-press robot, or, where that is not possible, with a browser timing test, and the camera's frame rate is stored with every minute. Participants are told that BlinkTwin must not be used to decide whether they are fit to drive. Bystanders who may appear on the webcam at work, and the employer where relevant, are informed and asked for consent; no video is stored in the field phase.

\paragraph*{Laboratory phase} A subset takes part in a supervised session of up to 26~h of continuous wakefulness after at least 7 days of a stable sleep schedule verified by actigraphy, with simultaneous EEG, EOG and video. During the session the camera records desk work (typing, reading, phone use), not only the PVT or resting EEG. Additional exclusions apply (epilepsy, bipolar disorder, cardiovascular disease), caffeine and alcohol are not allowed from 24~h before, a physician is on call, and participants are taken home by arranged transport. Long eye closures and microsleeps are scored from EEG, EOG and video by two blinded raters using the BERN criteria~\cite{hertig2020microsleep}, and inter-rater agreement is reported. The number of participants is set by events, not by people: recruitment continues until at least 100 scored long eye closures are observed, which gives a 95\% confidence interval half-width of about 0.1 for a sensitivity near 0.7; analysis is per event with clustering by participant. Identifiable video is stored encrypted on institutional servers, used only for scoring and deleted after the study.

\paragraph*{Outcomes} Two co-primary outcomes: (1) the absolute error of predicting each probe rating, made from all earlier events (leave-future-out replay of the full 10-day log), with rating noise included in the interval; and (2) the absolute error of predicting the median reaction time of the next PVT-B, an objective outcome. Secondary outcomes: the error of predicting check-in ratings; coverage of the 90\% prediction interval, the probability integral transform and the continuous ranked probability score; the Brier score and AUROC for KSS $\geq 7$ within the next hour, in the laboratory phase where such ratings are common; sensitivity and positive predictive value of the alerts against scored events (laboratory) and alerts per hour with participant-labelled false alarms (field); agreement of the camera's blink measures with EOG; the false-flag rate of the reactions-versus-ratings check; the error of the predicted next TLX rating (exploratory); usability (System Usability Scale~\cite{brooke1996sus}); and adherence.

\paragraph*{Hypotheses and analysis} Methods are those of Section~\ref{sec:simulation}: the prior alone, check-ins only, the consider-covariance comparator and BlinkTwin. H1: on probe ratings, BlinkTwin is not worse than check-ins only, with a non-inferiority margin of 0.05 KSS points on the per-person mean absolute error (31\% of the simulated benefit of check-ins over the prior, 0.16 points), tested with a one-sided 97.5\% confidence bound and reported following the CONSORT extension for non-inferiority~\cite{piaggio2012reporting}. H2: on probe ratings, BlinkTwin is better than check-ins only, which is the test of the passive signals. H3: BlinkTwin predicts the next PVT-B better than the prior alone. H1 is tested first; H2 and H3 are tested at a Bonferroni-corrected $\alpha$ of 0.025 each. Absolute errors are also analysed with linear mixed models with method as a fixed effect and random intercepts for participant and day, as a sensitivity analysis. No parameter is tuned on the study data.

\paragraph*{Sample size} We estimated power by simulating the analysis that will be run, for check-in ratings: 1000 virtual people (the workday population of Section~\ref{sec:simulation}) each had 10 workdays replayed as one continuous event log, so that the offset and anchors carried over between days, with four check-ins a day each followed by a PVT-B; the prediction of each observed, noisy and rounded rating used all earlier events; and studies of $n$ people were drawn from the pool 5000 times. The mean error was 1.01 for the prior alone, 0.85 for check-ins only and 0.84 for BlinkTwin; per person, BlinkTwin minus check-ins averaged $-0.010$ (SD 0.036) and BlinkTwin minus the consider-covariance comparator $-0.014$ (SD 0.031). With a mean-reverting offset the difference from check-ins grew to $-0.028$ (24-h time constant) and $-0.047$ (72~h). Replaying each day separately instead, as in our earlier analysis, reversed the sign ($+0.019$), because the anchors were then relearnt every morning. Non-inferiority at a 0.05 margin had power above 0.99 from 20 people. Superiority of BlinkTwin over check-ins only had power 0.41 with 40, 0.68 with 80, 0.76 with 100 and 0.93 with 150 people (two-sided $\alpha = 0.05$). The co-primary outcomes on probe ratings and the next PVT-B, and the Bonferroni-corrected $\alpha$, have not yet been simulated; they will be before registration. Provisionally we plan to analyse 120 participants and to recruit 135. The generator keeps each person's offset constant across days, which the weak persistence in ds004902 contradicts, so these figures are likely optimistic for BlinkTwin.
